\subsection{同一域的赋值环的包含}

命题 1. 设 K 是域，A 是 K 的一个赋值环。那么：

\begin{enumerate}
\def\labelenumi{(\alph{enumi})}
\item
  每个环 \(\mathbf{B}\) ，只要它满足 \(\mathbf{A} \subset \mathbf{B} \subset \mathbf{K}\) ，就是 \(\mathbf{K}\) 的赋值环；
\item
  这种环的极大理想 \(\mathfrak{m}(\mathbf{B})\) 含于 \(\mathbf{A}\)，且它是 \(\mathbf{A}\) 的素理想；
\item
映射 \(\mathfrak{p} \mapsto \mathbf{A}_{\mathfrak{p}}\) 是由 \(\mathbf{A}\) 的素理想集合到环 \(\mathbf{B}\) （满足 \(\mathbf{A} \subset \mathbf{B} \subset \mathbf{K}\) ）的集合上的一个递减双射；其逆双射是映射 \(\mathbf{B} \mapsto \mathfrak{m}(\mathbf{B})\) 。
\end{enumerate}

设 B 是满足 \(A \subset B \subset K\) 的环且 \(x \in K - B\) ，则 \(x \in K - A\) ，于是 \(x^{-1} \in \mathfrak{m}(A) \subset B\) ，这同时证明了 B 是 K 的赋值环以及 \(\mathfrak{m}(B) \subset \mathfrak{m}(A)\) ；由于 \(\mathfrak{m}(B) = \mathfrak{m}(B) \cap A\) 是 A 的素理想，我们已证明了 (a)

与 (b)。此外，\(\mathbf{A}_{\mathfrak{m}(\mathbf{B})} \subset \mathbf{B}\) ；反之，若 \(x \in \mathbf{B} - \mathbf{A}\) ，则 \(x^{-1} \in \mathbf{A}\) 且 \(x^{-1} \notin \mathfrak{m}(\mathbf{B})\) ，因此 \(x \in \mathbf{A}_{\mathfrak{m}(\mathbf{B})}\) ；于是 \(\mathbf{A}_{\mathfrak{m}(\mathbf{B})} = \mathbf{B}\) 。最后设 \(\mathfrak{p}\) 是 \(\mathbf{A}\) 的素理想；记 \(\mathbf{B} = \mathbf{A}_{\mathfrak{p}}\) ；则 \(\mathfrak{m}(\mathbf{B}) \cap \mathbf{A} = \mathfrak{p}\) （第二章 §2 第 5 节命题 11），且由 (b) 有 \(\mathfrak{m}(\mathbf{B}) \subset \mathbf{A}\) ；因此 \(\mathfrak{m}(\mathbf{B}) = \mathfrak{p}\) ，这表明命题中的映射 \(\mathfrak{p} \mapsto \mathbf{A}_{\mathfrak{p}}\) 与 \(\mathbf{B} \mapsto \mathfrak{m}(\mathbf{B})\) 是互逆的双射。

推论. K 中含 A 的子环的集合按包含关系是全序的。

A 的素理想集合按包含关系是全序的（§ 1 第 2 节定理 1 (e)），且映射 \(\mathfrak{p} \mapsto \mathrm{A}_{\mathfrak{p}}\) 反转包含关系。

命题 2. 设 K 是域，B 是 K 的一个赋值环，\(h_B\) 是与 B 相伴的 K 的位（取值于 \(\kappa(B)\) ）。那么映射 \(A \mapsto h_B(A)\) 定义了由含于 B 的 K 的赋值环组成的集 \(\mathfrak{A}\) 到 \(\mathfrak{A}'\) （\(\kappa(B)\) 的赋值环组成的集）上的一个双射。

若 \(\mathbf{A} \in \mathfrak{A}\) ，则 \(h_{\mathrm{B}}(\mathbf{A}) \in \mathfrak{A}'\) ：因为若 \(x' = h_{\mathrm{B}}(x)\) （其中 \(x \in \mathbf{B}\) ）是 \(\kappa(\mathbf{B}) - h_{\mathrm{B}}(\mathbf{A})\) 的元素，则 \(x \notin \mathbf{A}\) ，于是 \(x^{-1} \in \mathbf{A}\) ，且 \(h_{\mathrm{B}}(x)^{-1} \in h_{\mathrm{B}}(\mathbf{A})\) 。另一方面，对 \(\mathbf{A} \in \mathfrak{A}\) 有 \(\mathbf{A} \supset \mathfrak{m}(\mathbf{B})\) （命题 1 (b)），因此映射 \(\mathbf{A} \mapsto h_{\mathrm{B}}(\mathbf{A})\) 是单射。最后，设 \(\mathbf{A}' \in \mathfrak{A}'\) ，\(\mathbf{A} = \bar{h}_{\mathrm{B}}^{-1}(\mathbf{A}') \subset \mathbf{B}\) ；我们将证明 \(\mathbf{A} \in \mathfrak{A}\) ，从而完成证明；若 \(x \in \mathbf{K} - \mathbf{A}\) ，则或者 \(x \notin \mathbf{B}\) ，或者 \(x \in \mathbf{B}\) ；若 \(x \notin \mathbf{B}\) ，则 \(x^{-1} \in \mathfrak{m}(\mathbf{B}) \subset \mathbf{A}\) ；若 \(x \in \mathbf{B}\) ，则 \(h_{\mathrm{B}}(x) \in \kappa(\mathbf{B})\) 且 \(h_{\mathrm{B}}(x) \notin \mathbf{A}'\) ，于是 \(h_{\mathrm{B}}(x^{-1}) \in \mathbf{A}'\) ，我们再次得出 \(x^{-1} \in \mathbf{A}\) ；因此 \(\mathbf{A} \in \mathfrak{A}\) 。

推论. 设 A 与 B 是 K 的两个赋值环，其中 \(A \subset B\) ；设 \(A' = h_{B}(A)\) ，它是 \(\kappa(B)\) 的一个赋值环。\(\kappa(A')\) 是 \(A'\) 的剩余域，它典范地同构于 A 的剩余域 \(\kappa(A)\) ，且与 A 相伴的位 \(h_{A}\) 是 \(h_{A'} \circ h_{B}\) （与 \(A'\) 和 B 相伴的位的复合）。

由于局部环 A' 是局部环 A 的一个商，它们的剩余域典范地同构，且等式 \(h_{\mathrm{A}}(x) = h_{\mathrm{A}'}(h_{\mathrm{B}}(x))\) 对 \(x \in A\) 成立。另一方面，若 \(x \in B - A\) ，则 \(h_{\mathrm{B}}(x) \notin A'\) ，等式两边都等于 \(\infty\) ；若 \(x \in K - B\) ，情形相同。

注. 反之，设 \(f\) 是 \(\mathbf{K}\) 的一个取值于 \(\mathbf{K}'\) 的位，\(f'\) 是 \(\mathbf{K}'\) 的一个取值于 \(\mathbf{K}''\) 的位。那么 \(f' \circ f\) 是 \(\mathbf{K}\) 的一个位，它的环含于位 \(f\) 的环。

